\documentclass[11pt]{article}
\usepackage[margin=1in]{geometry}
\usepackage{amsmath,amssymb,amsthm}
\newtheorem{theorem}{Theorem}
\title{A disproof of Conjecture 00000000961}
\author{earthking11}
\date{October 6, 2026}
\begin{document}
\maketitle
\section*{Third cumulants}
The conjecture assigns the coefficient $\sqrt2$ to the ratio of the third
free cumulant and the third classical cumulant.  At order three, however, the
two moment--cumulant formulae coincide:
\[
 \kappa_3=m_3-3m_2m_1+2m_1^3,
 \qquad
 r_3=m_3-3m_2m_1+2m_1^3.
\]
The distinction between all set partitions and noncrossing partitions first
affects the cumulant formula at order four, not order three.

\begin{theorem}
The claimed universal ratio $\sqrt2$ is false.
\end{theorem}
\begin{proof}
Let $X=-1$ with probability $2/3$ and $X=2$ with probability $1/3$.  Then
\[
 m_1=0,\qquad m_2=2,\qquad m_3=2.
\]
Substitution gives $\kappa_3=r_3=2$.  Their ratio is therefore $1$.  Equivalently,
the necessary identity $r_3^2=2\kappa_3^2$ would read $4=8$.
\end{proof}

\section*{Lean 4 certificate}
The project in \texttt{lean4/} checks the weighted moments as exact integer
identities after multiplication by $3$, implements both order-three formulas,
and proves $4\ne8$.  The certificate is core Lean and has no unproved
placeholder or added axiom.
\end{document}
